\section{Results}
\label{sec:results}

% ── 3.1 Corpus coverage ──────────────────────────────────────────────────────
\subsection{Corpus coverage and index statistics}

Full ingestion of all 331 documentation sources (Table~\ref{tab:sources}),
including the optional Powder Crystallography textbook (\texttt{--book}),
completes in approximately 30 minutes on a modern laptop CPU (single core),
producing a $\sim$65\,MB ChromaDB collection of 5{,}838 embedded
text chunks with a median chunk length of approximately 810 characters.
The default configuration (without \texttt{--book}) indexes 152 sources
and produces $\approx$5{,}060 chunks in under 10 minutes.

The Programmer's Guide contributes the largest single block of content
(2{,}740 chunks) because the ReadTheDocs API documentation pages contain
dense inline docstrings, parameter tables, and code examples across
23 module-level pages.
Tutorial sections tend to be long (500--1{,}200 chars after splitting)
while many help-manual entries are short definitional paragraphs
($<$200 chars).
Short chunks are retained without further filtering because they
frequently contain precise parameter definitions or menu-path
specifications that are highly relevant to targeted queries.

% ── 3.2 Retrieval quality ────────────────────────────────────────────────────
\subsection{Retrieval quality}
\label{sec:retrieval-eval}

We constructed a held-out evaluation set of 40 questions spanning the
major GSAS-II workflow categories: Rietveld refinement (10),
sequential refinement (8), structure solution (6), instrument
parameter calibration (6), scripting (5), and data export (5).
Each question was paired with one or more ground-truth source
URLs selected from the indexed documentation corpus that contain
the correct answer; questions were verified against the tutorials
and help pages to confirm those answers appear in the corpus.

For each question we assessed whether the ground-truth source page
appeared among the top-$k$ retrieved results for $k \in \{1, 3, 6\}$
(Recall@$k$) and the rank of the first correct hit (Mean Reciprocal
Rank, MRR).
The evaluation set is available in the repository at
\texttt{evaluation/questions.json}.
Table~\ref{tab:retrieval} reports results measured against
the full 5{,}838-chunk index (including the \textit{Powder Diffraction
Crystallography} textbook).

\begin{table}[h]
  \centering
  \caption{Retrieval performance on the 40-question evaluation set
           using the \texttt{BAAI/bge-base-en-v1.5} embedding model
           and the full index (5{,}838 chunks including the textbook).
           $k$ is the number of retrieved chunks inspected.
           Numbers are proportions (hits / questions in category).}
  \label{tab:retrieval}
  \begin{tabular}{lccc}
    \toprule
    & Recall@1 & Recall@3 & Recall@6 \\
    \midrule
    All categories ($n=40$)      & 0.50 & 0.83 & 0.98 \\
    Rietveld ($n=10$)            & 0.40 & 0.70 & 1.00 \\
    Sequential refine.\ ($n=8$)  & 0.38 & 0.75 & 1.00 \\
    Structure solution ($n=6$)   & 0.33 & 1.00 & 1.00 \\
    Calibration ($n=6$)          & 1.00 & 1.00 & 1.00 \\
    Scripting ($n=5$)            & 0.40 & 0.80 & 0.80 \\
    Data export ($n=5$)          & 0.60 & 0.80 & 1.00 \\
    \midrule
    MRR (all)           & \multicolumn{3}{c}{0.66} \\
    \bottomrule
  \end{tabular}
\end{table}

Recall@6 of 0.98 (39 of 40 questions) indicates that for almost all user
questions, at least one of the six retrieved chunks contains the relevant
answer, regardless of LLM synthesis quality.
The Calibration category achieves perfect Recall@1 (1.00), with the
correct source page ranked first for every calibration question.
The lower Recall@1 overall (0.50) reflects competition from the tutorials
index page (\texttt{tutorials.html}) --- which summarises all tutorials in
one place and frequently ranks above individual tutorial pages for
broad queries --- and from semantically similar sub-workflows that share
vocabulary (\emph{e.g.}\ Monte Carlo simulated annealing versus
RMCProfile for structure solution).
Retrieval is performed by over-fetching the top-30 candidates and
applying a URL-diversity step that retains only the highest-scoring chunk
per unique source page before selecting the final six; this prevents any
single high-chunk-count page from dominating the result set.
The one remaining miss (Scripting Q33: ``read back Rwp and lattice
parameters from a script'') arises because the vocabulary of those
quantities appears throughout the tutorial corpus; the scripting context
is overwhelmed by frequency.
The Scripting category Recall@6 (0.80) reflects this single structural
failure rather than a general weakness in scripting retrieval.

% ── 3.3 Answer quality ───────────────────────────────────────────────────────
\subsection{Answer quality and inline citations}

Figure~\ref{fig:screenshot} shows a representative query and response
using the Llama 3 8B backend.
The answer correctly identifies the \textit{Calculate / Refine} menu
path, the expected Rwp convergence value, and the Le Bail
intensity initialisation step, with each claim supported by an
inline citation.
In a manual review of 20 sampled answers by two expert GSAS-II users,
all factual assertions that carried an inline citation were judged
to be accurate and traceable to the cited section.
Unfounded assertions --- present in 3 of 20 responses --- all
occurred in clauses not assigned a citation number, consistent
with the model using its parametric knowledge rather than the
retrieved context for those clauses.

The citation injection post-processor (Section~\ref{sec:citations})
increased the mean number of inline citation markers per answer from
1.2 (model-only, Qwen2.5-3B backend) to 5.7 (all retrieved chunks
assigned at least one citation), without introducing false attributions
in the 20 reviewed responses.
Citation numbers are renumbered by first appearance order so that
$[1]$ always precedes $[2]$ in the answer text, matching the
expectations of numbered-citation academic style.

Figure~\ref{fig:webapp} shows the web UI displaying inline citations
rendered as clickable superscript links above a journal-style
numbered bibliography.
Figure~\ref{fig:gui} illustrates the wxPython desktop dialog responding
to a mathematical query; when internet access is available, the
\textit{Query-GSAS} GUI uses MathJax to render \LaTeX{} equation markup
produced by the LLM, improving readability for formula-heavy answers
such as the Debye scattering equation.

% ── 3.4 Inference latency ────────────────────────────────────────────────────
\subsection{Inference latency}

End-to-end response latency (retrieval + LLM generation, median over 20 queries)
is approximately 8\,s for Llama 3 8B via Ollama on a modern laptop CPU,
and under 3\,s with GPU acceleration, making the system practical for
interactive use during active data analysis.
The retrieval-only backend completes in under 0.5\,s and serves users
who prefer to read raw documentation excerpts directly.
